\section{RewardBench：奖励模型的评估}

\subsection{为什么需要评估奖励模型？}

Lambert 反复强调一个观点：\textbf{工业界声称奖励模型是 RLHF 成功的关键，但我们却缺乏系统评估奖励模型的工具。}

\begin{figure}[H]
\centering
\includegraphics[width=0.95\textwidth]{slides-images/slide-035.jpg}
\caption{Life after DPO 的两大方向：更好的对齐评估 + 如何改进 DPO 模型\protect\footnotemark}
\end{figure}
\footnotetext{Slides 第36页。}

在标准的 RLHF 反馈循环中：策略模型生成回复 $\rightarrow$ 奖励模型评分 $\rightarrow$ 策略更新。但现有的评估工具（如 AlpacaEval、MT-Bench、Chatbot Arena）都是评估最终策略的——\textbf{没有工具直接评估奖励模型本身}。

\begin{figure}[H]
\centering
\includegraphics[width=0.95\textwidth]{slides-images/slide-037.jpg}
\caption{RLHF 反馈循环：评估工具缺失对奖励模型内部环节的洞察\protect\footnotemark}
\end{figure}
\footnotetext{Slides 第38页。}

\begin{importantbox}{为什么评估奖励模型如此重要？}
训练模型时，你不可能等 Chatbot Arena 一个月后返回结果来判断训练是否有效。你需要\textbf{本地评估工具}——能在训练过程中快速给出信号，告诉你技术改进是否真正有效。评估奖励模型是理解和改进整个 RLHF 流程的关键环节。
\end{importantbox}

\subsection{奖励模型的训练方式}

奖励模型的训练有几个独特的特点：

\begin{figure}[H]
\centering
\includegraphics[width=0.95\textwidth]{slides-images/slide-040.jpg}
\caption{奖励模型的训练损失函数\protect\footnotemark}
\end{figure}
\footnotetext{Slides 第41页。}

训练奖励模型需要\textbf{成对偏好数据}：一个 prompt 配两个 completion（chosen 和 rejected）。模型同时处理两个输入，输出两个标量值，通过损失函数拉大它们之间的距离：

$$L_{\mathrm{PM}} = \log(1 + e^{r_{\mathrm{rejected}} - r_{\mathrm{chosen}}})$$

\begin{knowledgebox}{奖励模型训练的特殊性}
\begin{itemize}
    \item 通常\textbf{只训练 1 个 epoch}——多训会过拟合
    \item 评估时与标注者的\textbf{一致性只有 65--75\%}
    \item 可以添加 margin loss、使用模型集成等技巧，但效果提升有限
    \item 低一致性可能不完全是 bug——人的偏好本身就存在分歧，模型不应该过于"狭隘"
\end{itemize}
\end{knowledgebox}

\subsection{RewardBench 的设计与结果}

RewardBench 的设计很直接：收集一组 prompt，每个 prompt 手动创建 chosen 和 rejected 回复，然后测试奖励模型是否能正确识别。

\begin{figure}[H]
\centering
\includegraphics[width=0.95\textwidth]{slides-images/slide-045.jpg}
\caption{RewardBench 发布时（2024年3月）的排行榜\protect\footnotemark}
\end{figure}
\footnotetext{Slides 第46页。}

排行榜中包含多种类型的模型：
\begin{itemize}
    \item \textbf{序列分类器}（Sequence Classifier）：传统的奖励模型训练方式
    \item \textbf{DPO 模型}：可以利用隐含的奖励函数进行评估
    \item \textbf{生成模型}（Generative Model / LLM-as-a-judge）：直接用语言模型判断哪个回答更好
\end{itemize}

\begin{figure}[H]
\centering
\includegraphics[width=0.95\textwidth]{slides-images/slide-050.jpg}
\caption{RewardBench 更新版（2024年5月）：两个月内，原来排名第5的模型跌至第31位\protect\footnotemark}
\end{figure}
\footnotetext{Slides 第51页。}

一个关键发现是：\textbf{GPT-4 和 GPT-4o 在作为 LLM-as-a-judge 时，表现并不如 Cohere 专门训练的奖励模型}。这说明通用能力强的模型不一定是好的奖励模型——专门训练的小模型在这个任务上可能更有效。

\subsection{Chat Hard 子集：真正的难题}

\begin{figure}[H]
\centering
\includegraphics[width=0.95\textwidth]{slides-images/slide-052.jpg}
\caption{Chat Hard 数据示例：要求模型区分关于"星星"的隐喻和关于"月亮"的隐喻\protect\footnotemark}
\end{figure}
\footnotetext{Slides 第53页。}

RewardBench 中最有价值的子集是 \textbf{Chat Hard}——专门设计的"陷阱题"。例如：给定提示"用星星写一个比喻"，chosen 回复是关于星星的比喻，rejected 回复则是关于月亮的比喻。由于星星和月亮在语义空间中高度关联，这对语言模型来说是一个非常困难的判断任务。这个子集是 RewardBench 中唯一尚未被"饱和"的部分。

\subsection{安全性评估的洞察}

\begin{figure}[H]
\centering
\includegraphics[width=0.95\textwidth]{slides-images/slide-054.jpg}
\caption{安全性子集的结果：模型可以分为"全部拒绝"和"全部回答"两个极端\protect\footnotemark}
\end{figure}
\footnotetext{Slides 第55页。}

安全性评估中观察到的模式完全符合预期：有些模型（过度安全型）会拒绝所有请求，有些模型（"uncensored"型）会回答所有请求。理想的模型应该能够\textbf{区分}真正应该拒绝的请求和只是看起来敏感但实际上应该回答的请求。

\subsection{DPO 隐含奖励模型的局限}

\begin{figure}[H]
\centering
\includegraphics[width=0.95\textwidth]{slides-images/slide-055.jpg}
\caption{DPO 论文中的奖励定义（公式3和公式4）：需要参考模型才能正确计算\protect\footnotemark}
\end{figure}
\footnotetext{Slides 第56页。}

DPO 论文中定义的隐含奖励为：

$$r(x, y) = \beta \log \frac{\pi(y|x)}{\pi_{\mathrm{ref}}(y|x)} + \beta \log Z(x)$$

使用 DPO 模型作为奖励模型时，需要同时拥有训练后的策略和参考模型（训练前的检查点）。但实际上，\textbf{大多数人发布 DPO 模型时不会发布参考模型}。

\begin{warningbox}{没有参考模型的 DPO 无法作为好的奖励模型}
如果去掉参考模型（即划掉公式中的 $\pi_{\mathrm{ref}}$ 项），DPO 模型在 RewardBench 上的表现会\textbf{大幅下降}。这是因为参考模型在数学推导中起到关键的正则化作用。这意味着：要把 DPO 模型当作奖励模型使用，你必须同时拥有参考模型——但这在实践中很难要求所有模型发布者都做到。
\end{warningbox}

\subsection{本章小结}

RewardBench 填补了对齐研究中奖励模型评估的空白。核心发现包括：（1）专门训练的奖励模型优于通用 LLM-as-a-judge；（2）Chat Hard 子集是衡量奖励模型能力的关键指标；（3）DPO 模型作为奖励模型有数学上的局限性——需要参考模型才能正常工作；（4）奖励模型的评估领域仍在快速发展，两个月内排行榜就会发生巨大变化。

%% ========== Part 5: DPO vs PPO ==========

